\chapter{Methods}
\label{chap:method}
\section{Notation and setup}
\label{sec:notation}
In this section we introduce the notation and the setup used throughout this thesis. 
Let $i=1, \ldots, N$ index the individuals in the study and let $t=0, \ldots, T$ 
index the time points at which measurements are taken. We denote the treatment 
variable at time $t$ for individual $i$ as $A_{it}$, which can take on values from a 
predefined set of treatment options. In this thesis, 
we assume a binary treatment, $A_{it} \in \{0,1\}$. We adopt the convention that treatment 
is absent at the beginning of the study, $t=0$ and therefore $A_{i0} = 0$ 
for all individuals $i$. 
The outcome variable at 
time $t$ for individual $i$ is 
denoted as $Y_{i,t+1}$ to emphasize that the outcome is measured after the treatment at time $t$.
We also consider a set of observed covariates $X_{it}$, 
which may include both time-varying 
and time-invariant variables. The history of treatment and covariates up to time $t$ for 
individual $i$ is represented as
$\bar{A}_{it} = (A_{i1}, A_{i2}, \ldots, A_{it})$ and $\bar{X}_{it} = (X_{i1}, X_{i2}, \ldots, X_{it})$, 
respectively. We denote the potential outcome at time $t$ for individual $i$ under a specific treatment
regime $\bar{a}_t = (a_1, a_2, \ldots, a_t)$ as $Y_{i,t+1}(\bar{a}_t)$. Finally, we use $U_i$ 
to represent \emph{structural} unmeasured confounders that are time-invariant for individual $i$. 

We also define truncated versions of the treatment and covariate histories from \(k\) 
steps prior up to time \(t\) as \(\underbar{A}_{i,t-k}\) and \(\underbar{X}_{i,t-k}\), respectively. 
The outcome setup we are focusing on is a single outcome measured at the end of 
the study, i.e., at time T+1, denoted as $Y_{i,T+1}$. This is the setup considered in
\cite{blackwellEffectPoliticalAdvertising2024} and it is a special case of the setup in
\cite{hattSequentialDeconfoundingCausal2024} who consider outcomes measured at each time point. This 
choice is both motivated by the specific application they consider 
(political advertising) and technical reasons related to the consistency of the IPTW-FE estimator.

Moreover, since the asymptotic properties and consistency of the IPTW-FE estimator relies on the 
truncated potnetial outcomes, we can define the truncated potential outcomes as
$Y_{i,T+1}(\underline{a}_{T-k}) \equiv Y_{i,T+1}(\overline{A}_{i,T-k-1},\, \underline{a}_{i,T-k})$. This 
definition represents the potential outcome of manipulating 
the last $k$ treatment values, while leaving the earlier treatment values as observed. In this way, 
the earlier treatments before the last $k$ time points can be thought of as part of the baseline covariates. 
In this section, we describe the two proposed methods 
for adjusting for unmeasured time-invariant confounding: the fixed-effect marginal 
structural model (FE-MSM) proposed by \citeA{blackwellEffectPoliticalAdvertising2024}
and the sequential Gaussian process latent variable model (SeqGPLVM) 
proposed by \citeA{hattSequentialDeconfoundingCausal2024} when it is used 
in conjunction with IPTW estimation of MSM. We first describe the common set of
assumptions required by of these methods, and then go 
into the details of the steps of each method. The causal target of interest for both methods 
is the average truncated potential outcome.  

FE-MSM consists of two stages: (i) a treatment model that 
estimates the propensity scores, conditioning on both observed 
covariates and individual fixed effects that capture the unmeasured time-invariant confounding, 
and (ii) a marginal structural model that uses the estimated propensity scores in stage one to adjust for confounding using IPTW.

SeqGPLVM in combination with MSM on the other hand, consists of three stages: 
(i) a latent variable model (SeqGPLVM) that estimates a substitute 
for the unmeasured confounder based on the observed treatment sequence, (ii) estimating 
the propensity scores using a treatment model that conditions on both observed 
covariates and the estimated unmeasured confounder substitute, and (iii) 
a marginal structural model that uses the estimated propensity scores in stage two 
to adjust for confounding using IPTW. It is worth mentioning that SeqGPLM can be used with any
other estimation methods as well, such as g-computation or doubly robust methods, 
but we focus on IPTW estimation of MSM for comparability with FE-MSM.

\section{Common assumptions}
\label{sec:assumptions}
The common assumptions for both FE-MSM and SeqGPLVM-MSM are stated below. 
The assumptions can be categorized into two groups: identification assumptions and 
modeling assumptions. Identification assumptions, 
are the assumptions which are needed to write target causal estimands such as 
$\mathbb{E}[Y_i(\underline{\mathbf{a}}_{T-k})]$ in terms of observed data. 
Based on the identification strategy and the estimator of choice, 
certain quantities must be modeled. 
For example, estimating a marginal structural model for the target causal estimand using 
inverse probability of treatment weighting (IPTW) requires modeling the 
propensity scores. Modeling assumptions are about the 
requirements placed on these models in order to have unbiased or consistent estimates. 
Throughout the text, 
we use C, I, and M to denote common, identification, and modeling 
assumptions, respectively.
% --- Common assumptions (C1, C2, ...) ---
% Assumes you already defined the assumpC environment (C-prefix) in the preamble.

\begin{idassumpC}[Consistency]\label{ass:C-consistency}
The observed outcome for each individual under their observed treatment history equals
the potential outcome under that treatment history. Formally, if $\bar{A}_{iT}=\bar{\mathbf{a}}_{iT}$, 
then $Y_{i}=Y_{i}(\bar{\mathbf{a}}_{iT})$. 
Moreover, $Y_{i} = Y_{i}(\underline{\mathbf{a}}_{i,T-k})$ if 
$\bar{\mathbf{A}}_{iT}=\left(\bar{\mathbf{A}}_{i,T-k-1},\underline{\mathbf{a}}_{i,T-k}\right)$.
\end{idassumpC}

\begin{comment}

\begin{assumpC}[Positivity]\label{ass:C-positivity}
\[
\Pr\!\left[ A_t = a_t \,\middle|\, 
\bar{A}_{t-1} = \bar{a}_{t-1},\,
\bar{X}_t = \bar{x}_t \right] > 0
\]

for all values $(\bar{a}_{t-1}, \bar{x}_t)$ with

\[
\Pr\!\left[ \bar{A}_{t-1} = \bar{a}_{t-1},\,
\bar{X}_t = \bar{x}_t \right] \neq 0
\]

in the population of interest, and for all $a_t \in \mathcal{A}_t$.
\end{assumpC}

\end{comment}

\begin{idassumpC}[No interference]\label{ass:C-no-interference}
The treatment assignment for one individual does not affect the potential outcomes of
other individuals. This assumption is also known as the no spillover assumption.
\end{idassumpC}

\begin{idassumpC}[Time-invariance of unmeasured confounding]\label{ass:C-timeinvariant-uc}
The unmeasured confounding is assumed to be
time-invariant. That is, $\mathbf{U}_{i,t} = \mathbf{U}_i$ for $t \in \{1,\dots,T\}$. We can then write 
conditional ignorability as 
\[
\{ Y_{i,t+1}(\bar{\mathbf{a}}_{t}),\, \bar{\mathbf{X}}_{i,t+1}(\bar{\mathbf{a}}_{t}) \}
\perp\!\!\!\perp
A_{it}
\mid
\bar{\mathbf{X}}_{it},\,
\bar{\mathbf{A}}_{i,t-1} = \bar{\mathbf{a}}_{t-1},\,
\mathbf{U}_i .
\]
Note that this assumption is about the true underlying data-generating process and cannot be  
used for identification directly since $\mathbf{U}_i$ is unmeasured. 
\end{idassumpC}

\begin{idassumpC}[Positivity]\label{ass:C-positivity}
For all treatment times $t = T - k + 1, \ldots, T$, 
and all histories $(\bar{\mathbf{X}}_t, \bar{\mathbf{A}}_{t-1}, \boldsymbol{\eta}_i)$ 
with positive probability, every treatment level has positive probability:
\[
\mathbb{P}(A_t = a_t \mid \bar{\mathbf{X}}_t = \bar{\mathbf{x}}_t,\ \bar{\mathbf{A}}_{t-1} = 
\bar{\mathbf{a}}_{t-1}, \boldsymbol{\eta}_i) > 0 \quad \text{for all } a_t \in \mathcal{A}_t
\]
Where $\boldsymbol{\eta}_i$ represents the unit-specific fixed effect in the FE-MSM model 
or the confounder substitute in the SeqGPLVM model. This assumption  particularly
is necessary for the identification of $\mathbb{E}[Y_i(\underline{\mathbf{a}}_{T-k})]$.
\end{idassumpC}



\input{sections/method/first_stage.tex}
\input{sections/method/outcome.tex}
\section{Comparison of FE-MSM and SeqGPLVM}
\label{sec:comparison}
\paragraph{FE-MSM}
The Fixed-Effects MSM (FE-MSM) addresses time-invariant unmeasured confounding by incorporating unit-specific fixed effects directly into the propensity score model.

Unit-Specific Sequential Ignorability: This method relies on the assumption that treatment is randomized conditional on the observed history and a time-constant unit-specific effect $\alpha_i$.
Estimation Procedure: We obtain estimates for the propensity parameters $\beta$ and the unit effects $\alpha_i$ via **Conditional Maximum Likelihood**. This approach is justified under **large-$N$, large-$T$ asymptotics, where the bias from noisy fixed-effect estimation (the incidental parameters problem) fades over time.
Weighted Least Squares: Once the fixed-effect weights are constructed, the MSM parameters $\gamma$ are estimated by solving a weighted estimating equation, which reduces to **Weighted Least Squares (WLS)** for our linear specification. We use a **robust (sandwich) variance estimator** to obtain valid Wald confidence intervals.

\paragraph{SeqGPLVM}

In contrast to the point-estimate approach of FE-MSM, the SeqGPLVM captures unmeasured confounding by inferring a posterior distribution $q(Z)$ over a latent substitute $Z_i$.

Handling Circular Dependency: To avoid the "circular dependency" created by using $Z_i$ as a covariate in a separate propensity model, we use the estimated propensities from the SeqGPLVM **directly** to build the IPTW weights.
Multiple Imputation (MI)-Style Procedure: To propagate the latent uncertainty from the GP-based treatment model, we draw $S$ samples $Z^{(s)} \sim q(Z)$. For each draw, we:
    1.  Compute draw-specific propensities and weights $W^{(s)}$.
    2.  Fit the weighted MSM to obtain a draw-specific estimate $\hat{\gamma}^{(s)}$ and robust covariance $V^{(s)}$.
Rubin’s Rules for Pooling: The final point estimate is the average across draws, $\bar{\gamma} = S^{-1}\sum \hat{\gamma}^{(s)}$. The total variance is calculated using **Rubin’s within/between decomposition**, $T = W + (1 + 1/S)B$, which explicitly accounts for the variability induced by the uncertainty in the inferred latent confounder.



Although the seqgplvm can be used in the settings in which the outcome is observed at every time point, 
since the asymptotic properties and consistency of their IPTW-FE estimator 
relies on having less number of N than T (assumption )
Based on the consistency assumption \ref{ass:C-consistency}, we can define 
truncated potential outcomes 

\begin{comment}

\begin{figure}[ht]
\centering
\resizebox{\textwidth}{!}{
\begin{tikzpicture}[
    font=\small,
    % ADD THIS LINE:
    ampersand replacement=\&, 
    arr/.style={-Latex, thick},
    box/.style={
        draw, rounded corners, align=left,
        inner xsep=5mm, inner ysep=3.5mm,
        text width=0.42\linewidth % Changed from 0.44 to 0.42
    },
    header/.style={
        draw, rounded corners, align=center,
        inner xsep=5mm, inner ysep=2.5mm,
        text width=0.42\linewidth, % Changed from 0.44 to 0.42
        font=\bfseries
    },
    stage/.style={
        font=\bfseries,
        align=right,
        text width=0.20\linewidth, % Increased width
        anchor=east                % Added to fix alignment
    },
    midnote/.style={
        draw, rounded corners, fill=gray!10,
        align=center, inner xsep=2.5mm, inner ysep=2.5mm,
        text width=0.16\linewidth
    }
]

% --- Matrix: All '&' replaced with '\&' ---
\matrix (m) [
    matrix of nodes,
    nodes in empty cells,
    row sep=10mm,
    column sep=12mm,
    nodes={box, anchor=north}
] {
    |[header]| FE--MSM \& |[header]| SeqGPLVM \\
    {\textbf{MSM / estimand.}\\
    Choose $k$ and specify an MSM for $Y_i(d_{T-k:T})$.}
    \&
    {\textbf{MSM / estimand.}\\
    Choose $k$ and specify an MSM for $Y_i(d_{T-k:T})$.}\\
    {\begin{minipage}[t]{\linewidth}
        \begin{enumerate}\setlength\itemsep{1pt}\setlength\parskip{0pt}
        \item No interference
        \item Consistency
        \item Positivity
        \item \textbf{Unit-specific ignorability}
        \item \textbf{Truncated Marginal Structural Model}
        \end{enumerate}
    \end{minipage}}
    \&
    {\begin{minipage}[t]{\linewidth}
        \begin{enumerate}\setlength\itemsep{1pt}\setlength\parskip{0pt}
        \item No interference
        \item Consistency
        \item Positivity
        \item \textbf{Time-invariant unmeasured confounding}
        \item \textbf{Conditional Markov model}
        \item 
        \end{enumerate}
    \end{minipage}}\\
    {\textbf{Identification.}\\
    Given (claimed) ignorability, causal effects can be estimated\\
    by MSM/IPTW using propensities conditional on $(X, A, Z)$.}
    \&
    {\textbf{Identification.}\\
    Given (claimed) ignorability, causal effects can be estimated\\
    by MSM/IPTW using propensities conditional on $(X, A, Z)$. } \\
    {\textbf{Propensity estimation.}\\
    Specify and estimate $\pi_{it}(x_t, d_{t-1}; \beta, \alpha_i)$\\
    (FE MLE). }
    \&
    {\textbf{Stage 1: SeqGPLVM.}\\
    Fit GP+CMM latent factor model to infer $Z_i$ and/or compute\\
    $\hat{\pi}_{it}=\hat{p}(A_{it}\mid A_{i,t-1},X_{it},Z_i)$. } \\
    {\textbf{Construct weights.}\\
    IPTW / stabilized weights from $\hat{\pi}_{it}$\\
    (trim if needed). }
    \&
    {\textbf{Construct weights (this thesis).}\\
    Use $\hat{\pi}_{it}$ from SeqGPLVM directly to build IPTW weights.\\
    \emph{(Original: plug $\hat{Z}_i$ into an outcome model.)} } \\
    {\textbf{Outcome stage.}\\
    Fit MSM by weighted regression / estimating equations. }
    \&
    {\textbf{Outcome stage.}\\
    Fit MSM by weighted regression / estimating equations. } \\
};

% --- Row labels ---
\node[stage] at ($(m-2-1.west)+(-12mm,0)$) {Target};
\node[stage] at ($(m-3-1.west)+(-12mm,0)$) {Identification assumptions};
\node[stage] at ($(m-4-1.west)+(-12mm,0)$) {Estimation assumptions};
\node[stage] at ($(m-5-1.west)+(-12mm,0)$) {Treatment\\ model};
\node[stage] at ($(m-6-1.west)+(-12mm,0)$) {Weights};
\node[stage] at ($(m-7-1.west)+(-12mm,0)$) {Outcome};

% --- Vertical arrows ---
\foreach \r/\rp in {2/3,3/4,4/5,5/6,6/7} {
    \draw[arr] (m-\r-1.south) -- (m-\rp-1.north);
    \draw[arr] (m-\r-2.south) -- (m-\rp-2.north);
}

% --- Central "common output" callout ---
%\node[midnote] (pi) at ($(m-5-1.east)!0.5!(m-5-2.west)$) {Outputs\\$\hat{\pi}_{it}$};
%\draw[arr] (m-5-1.east) -- (pi.west);
%\draw[arr] (m-5-2.west) -- (pi.east);

\end{tikzpicture}
}
\caption{Stage-aligned comparison of FE--MSM and SeqGPLVM. The pipelines differ in how propensities are estimated (fixed effects vs. latent factor + GP), but both feed into IPTW and a weighted MSM fit.}
\label{fig:flow_compare_rows_clean}
\end{figure}
\end{comment}

